\scanpage{160}
\mspage{76}
suffit \inserted{en principe} de savoir exprimer les $V_{\Sigma',\Sigma''}$ et les morphismes de transition entre ceux-ci
\[
(259)\qquad V_{\Sigma',\Sigma''}\longrightarrow V_{\Sigma'_1,\Sigma''_1}
\]
pour
\[
(\Sigma',\Sigma'')\hookrightarrow(\Sigma'_1,\Sigma''_1),
\]
(quitte à appliquer ce formalisme, en remplaçant $X$ par $X\setminus\Sigma_0^{\unsure{'}}$ etc) -- on trouve alors les $V_{\Sigma',\Sigma''}^\Sigma$ comme le topos \deleted{\unsure{résiduel}} \inserted{ouvert de $V_{\Sigma',\Sigma''}^{\unsure{\Sigma}}$} relatif à l'immersion fermée
\[
V_{\Sigma'\cap\Sigma,\unsure{\Sigma'|\Sigma}}\longrightarrow V_{\Sigma',\Sigma''}.
\]
Mais cette ``réduction'' ne s'explicite pas aisément, au niveau des groupoïdes fondamentaux\ldots

Le fil directeur dans ce petit programme est celui-ci: on veut expliciter tous les topos qui sont associés à la filtration (comme ceux qu'on vient d'expliciter), et tout particulièrement leurs $\Pi_1$, en termes des topos \deleted{$\unsure{D_{d_0,\ldots,d_r}}$}
\[
V_{d_0,\ldots,d_r}\overset{\mathrm{déf}}=
V^{\deleted{\unsure{*}}}_{D^*_{d_0,\ldots,d_r^{\deleted{\unsure{*}}}},D_0}
\quad\text{et}
\]
\[
V^*_{d_0,\ldots,d_r}=V^*_{D^*_{d_0,\ldots,d_r},D_0}
\]
\margin{159}
